% ----------------------------------------------------------
\chapter{Resultados e Discussões}
% ----------------------------------------------------------

Os resultados apresentados no capítulo anterior demonstram que o protótipo de modelador de diagramas ER colaborativo entregou um sistema funcional que integra modelagem de dados, colaboração em tempo real e controle de acesso em uma única plataforma \textit{web}. Este capítulo discute esses resultados sob duas perspectivas: a comparação com ferramentas estabelecidas no mercado e a análise das dificuldades encontradas e das limitações do protótipo.

% ----------------------------------------------------------
\section{Comparação com Ferramentas Existentes}
% ----------------------------------------------------------

A fim de contextualizar o protótipo desenvolvido, a Tabela~\ref{tab:comparacao} apresenta uma comparação entre o protótipo e três ferramentas de modelagem ER disponíveis no mercado: Lucidchart, DBdiagram.io e MySQL Workbench.

\begin{table}[H]
\centering
\caption{Comparação entre o protótipo e ferramentas de modelagem ER}
\label{tab:comparacao}
\begin{tabular}{|l|c|c|c|c|}
\hline
\textbf{Funcionalidade} & \textbf{Protótipo} & \textbf{Lucidchart} & \textbf{DBdiagram.io} & \textbf{MySQL Workbench} \\
\hline
Modelagem ER              & Sim    & Sim      & Sim      & Sim \\
\hline
Colaboração em tempo real & Sim    & Pago     & Não      & Não \\
\hline
Exportação DDL (SQL)      & Sim    & Pago     & Sim      & Sim \\
\hline
Importação DDL (SQL)      & Sim    & Não      & Sim      & Sim \\
\hline
Baseado em \textit{web}   & Sim    & Sim      & Sim      & Não \\
\hline
Gratuito                  & Sim    & Parcial  & Parcial  & Sim \\
\hline
\end{tabular}
\begin{center}
    {\small Fonte: Elaborado pelo autor}
\end{center}
\end{table}

A análise comparativa evidencia que o protótipo desenvolvido distingue-se por reunir, em uma plataforma \textit{web} e de forma gratuita, as funcionalidades de colaboração em tempo real e de importação e exportação DDL. O Lucidchart, embora ofereça colaboração em tempo real, restringe esse recurso a planos pagos. O MySQL Workbench, apesar de robusto na modelagem e no suporte DDL, opera como uma aplicação \textit{desktop} e não oferece colaboração entre usuários. O DBdiagram.io, por sua vez, oferece suporte a DDL, mas carece da funcionalidade colaborativa. O protótipo preenche a intersecção dessas capacidades, alinhando-se ao diagnóstico de lacuna identificado na justificativa do trabalho: a ausência de uma ferramenta \textit{web}, acessível e colaborativa para a modelagem de diagramas ER.

% ----------------------------------------------------------
\section{Dificuldades e Lições Aprendidas}
% ----------------------------------------------------------

O desenvolvimento do protótipo apresentou desafios técnicos relevantes. A sincronização do estado do diagrama em um cenário de múltiplos colaboradores simultâneos constituiu a maior complexidade do projeto. A garantia de consistência --- propagando alterações de forma ordenada e prevenindo conflitos sem introduzir latência perceptível ao usuário --- exigiu iterações de projeto sobre a estratégia de \textit{locks} em memória e sobre a estrutura de mensagens STOMP. O modelo de \textit{locks} com \textit{timeout} automático demonstrou ser uma solução pragmática e suficiente para o escopo do protótipo.

A integração da biblioteca GoJS com o mecanismo de detecção de mudanças do Angular também demandou atenção. Alterações no modelo do diagrama originadas externamente, via WebSocket, precisaram ser aplicadas ao canvas da GoJS de maneira a não disparar ciclos redundantes de atualização, o que exigiu controle explícito das zonas de execução do Angular.

Por fim, a adoção da persistência poliglota --- PostgreSQL e MongoDB --- conferiu flexibilidade ao modelo de dados e desempenho adequado a cada tipo de operação. No entanto, essa decisão impôs a necessidade de gerenciar dois sistemas de banco de dados distintos, tanto no ambiente de desenvolvimento local quanto na configuração do Docker Compose para implantação, aumentando a complexidade operacional do projeto.

% ----------------------------------------------------------
\section{Limitações do Protótipo}
% ----------------------------------------------------------

Apesar dos resultados alcançados, o protótipo apresenta limitações que definem seu escopo em relação a uma solução de produção.

A primeira limitação refere-se à \textbf{ausência de um histórico de versões do diagrama}. Na versão atual, o estado do diagrama é substituído a cada atualização persistida, sem preservação de versões anteriores. Isso impede que os colaboradores desfaçam alterações em nível colaborativo ou inspecionem o histórico de mudanças do projeto, funcionalidade presente em ferramentas maduras de colaboração. A implementação de um mecanismo de versionamento, como o armazenamento de \textit{snapshots} incrementais, constitui uma evolução natural do protótipo.

A segunda limitação diz respeito à \textbf{licença comercial da biblioteca GoJS}. A GoJS exige a aquisição de uma licença para uso em ambiente de produção, disponibilizando apenas uma licença de avaliação para fins acadêmicos e de prototipagem. A adoção do protótipo em contexto comercial estaria sujeita a esses custos de licenciamento. Em trabalhos futuros, a substituição da GoJS por uma alternativa de código aberto com capacidades equivalentes seria um caminho viável para eliminar essa restrição.
